%% kapa: exam=1 round=22 year=2011 period=1 form=A subject=civil_law subject_ko=민법 seq=18 number=18
%% source: Q-Net 감정평가사 기출문제 (cst00309) artlSeq=5036790 file=extracted/1008/5036790/2011년 제22회 감정평가사 1차시험 기출문제/1교시.hwp
%% conversion: kordoc 4.24.51 -> tools/gen1.py (auto); review=proofread-sonnet; flags=text-verified-hwp
\begin{question}{18}
근저당권에 관한 설명으로 옳지 않은 것은? (다툼이 있으면 판례에 의함)
\begin{choices}
  \item 근저당권의 목적물이 양도된 후 피담보채무가 소멸한 경우, 근저당권설정자는 근저당권설정등기의 말소를 청구할 수 있다.
  \item 피담보채무를 확정시키는 근저당권설정자의 근저당권설정계약의 해제 또는 해지에 관한 권한은 근저당부동산의 소유권을 취득한 제3취득자도 원용할 수 있다.
  \item 근저당권자가 피담보채무의 불이행을 이유로 경매신청을 한 때에는 매수인이 매각대금을 완납한 때에 피담보채권은 확정된다.
  \item 근저당권의 피담보채권이 확정되기 전에 발생한 원본채권에 관하여 확정 후에 발생하는 이자나 지연손해금 채권은 채권최고액의 범위 내에서 근저당권에 의하여 담보된다.
  \item 확정된 피담보채권액이 채권최고액을 초과한 경우, 물상보증인이 채권최고액을 변제한 때에도 근저당권설정등기의 말소를 청구할 수 있다.
\end{choices}
\end{question}
\answer{3}
